% Per-task reference for Section 2. The main table (Table~\ref{tab:tasks}) carries the
% family, the starting model and the final metric; what it cannot hold is what each
% repository actually does, which quantity the agent can query while it works, and
% where the line between the two is drawn. Every number here is from the task's own
% instruction.md and the baseline row of Table~\ref{tab:grid}.
\section{The ten tasks}
\label{app:tasks}

Each task freezes one research repository and asks for the same thing: improve the training
algorithm that repository applies to its own model. What differs between them is the algorithm, the
asset the agent is given, and the pair of metrics either side of the evaluation boundary --- a cheap
one it may query as often as it likes during its four hours, and the one that decides its result,
computed afterwards by an evaluator it never sees. The two are related differently on different
tasks, and the relation matters when reading a column: on some the proxy is a subsample of the final
protocol, on others it is a different benchmark entirely.

\paragraph{OpenR1 --- supervised fine-tuning.}
Qwen2.5-Coder-1.5B-Instruct is fine-tuned on a decontaminated 8,005-row Python CodeForces
projection; the shipped recipe is completion-only supervised fine-tuning with the prompt tokens
masked. A candidate may select, reweight, pack, transform or synthesise training signal from those
rows, and may change the masking or the objective. The proxy is
\texttt{livecodebench\_public\_pass\_at\_1}; the final metric is the whole LiveCodeBench v6 release
slice, 175 problems under the benchmark's own sampling protocol --- ten samples per problem at
temperature $0.2$ and top-$p$ $0.95$, capped at 2,048 new tokens, scored as the mean over problems
of the fraction of samples that pass every official test.

\paragraph{RAGEN --- multi-turn agentic RL.}
A Qwen2.5-3B-Instruct policy is trained on Sokoban with multi-turn on-policy GRPO, generating its
own boards and trajectories online. Board construction, curriculum, rollout collection, reward
shaping and the update rule are all open; the scoring engine, the action decoding and the final
seeds are not. The proxy is a four-bank solve rate, the final metric a held-out 512-board solve
rate, and the two use different fixed environment-seed protocols.

\paragraph{OPD --- on-policy distillation.}
A 1.5B student is distilled from a mounted teacher by sampled-token on-policy distillation. The
proxy and the final metric are different benchmarks rather than two views of one: the proxy is
MATH-500 at four samples per question with a 12,288-token cap, and the final metric is AIME 2024 and
2025, 60 questions at 32 samples with a 31,744-token cap. MATH-500 is mounted during exploration;
the AIME inputs are not.

\paragraph{BTRM --- Bradley--Terry reward modelling.}
A scalar reward model is trained from a fixed Mistral-7B start on decontaminated UltraFeedback
preference pairs under a Bradley--Terry objective. The artifact must remain loadable as a scalar
reward model on that architecture, and any overlap between a training row and RewardBench
invalidates the run. Here the proxy is a strict subsample of the final: 512 of the 2,985 pairs are
visible, and the remaining 2,473 are held out until scoring.

\paragraph{DPO --- preference optimization.}
A merged Zephyr/Mistral-7B model is aligned with direct preference optimization. The final metric is
IFEval prompt-level strict accuracy over 413 held-out prompts.

\paragraph{DDPO --- diffusion RL.}
Stable Diffusion v1.5 is fine-tuned with on-policy DDPO and a LoRA adapter against a frozen CLIP
aesthetic reward. Prompt construction and sampling, reward shaping and normalization, auxiliary
losses, the update rule and the trainable parameters are all open; the reward assets and the final
prompt/latent stream are fixed, the latter mounted only at scoring. The proxy scores 64 generated
images, the final metric 256. CLIP alignment and mean pairwise image distance are reported alongside it.

\paragraph{NPO --- machine unlearning.}
Llama-3.2-1B-Instruct is unlearned on the TOFU \texttt{forget10} protocol with the official
OpenUnlearning NPO recipe. Final evaluation reports two pinned components, an extraction strength
that is better lower and a model utility that is better higher; the scalar we compare on is the
balanced score of Table~\ref{tab:tasks}, their harmonic mean after normalising each against the
training start and a retain-90 reference. During
exploration only the published anchor and the train-role projection are mounted --- the retain-90
anchor and the final-role data appear only in the separate score phase.

\paragraph{DiGress --- discrete graph diffusion.}
A discrete graph diffusion model is trained on QM9 without hydrogens. The fast metric is the
product of validity, uniqueness and novelty rates, higher better; the final metric is the upstream
test negative log-likelihood, lower better, and the repository's validation NLL is what connects
them. The real test split is mounted only at scoring.

\paragraph{Model Soup --- weight averaging.}
Seventy-two fixed CLIP ViT-B/32 checkpoints are combined into one model; the shipped construction is
a uniform mean. What the submitted code decides is which ingredients to use and with what
coefficients, and those coefficients may be negative or extrapolative. The proxy is 2,000
ImageNet-V2 images and the final metric the full 10,000.

\paragraph{OWL --- one-shot pruning.}
Unstructured sparsity is imposed on a dense OPT-6.7B in a single pass by activation-aware OWL/Wanda
pruning, with no fine-tuning in the shipped recipe. The hard artifact gate is decoder sparsity
within $[0.699, 0.701]$; within that window the pruning criterion, the search, the use of the
mounted C4 calibration shard, and any training a candidate cares to add are open. The proxy is
WikiText-2 validation perplexity and the final metric WikiText-2 test perplexity, the test text
mounted only at scoring.

\paragraph{}
Two of the ten ship a procedure that trains nothing: weight averaging combines checkpoints handed
over as data, and one-shot pruning removes weights from a released model in a single pass. Nothing
in the contract requires a submission to leave it that way, and on both tasks some did not --- the
submissions read in \S\ref{sec:reached} add a distillation fine-tune to the pruning pipeline and a
gradient-based search to the soup, and use the twelve hours accordingly. The two are kept because
the algorithmic question is as real in them as anywhere else: which checkpoints to combine and how,
which weights to remove and by what criterion.

% Z_task_contract.tex -- one task contract reproduced verbatim.
% The file is the agent-visible instruction.md of tasks/ragen_sokoban_grpo, unedited.
\section{One task contract in full}
\label{app:contract}

Every task carries the same contract, and about seventy per cent of its text is shared word for word
across the ten. Reproduced below is the whole of one of them, exactly as the agent receives it:
\texttt{instruction.md} for multi-turn agentic RL.

\begin{lstlisting}[style=contract]
# RAGEN on Sokoban

Improve the fixed Qwen2.5-3B-Instruct policy on the frozen Sokoban evaluation protocol. The shipped solution uses multi-turn on-policy GRPO and generates its training boards and trajectories online; that is the reference method rather than a mandatory objective.

You have up to four hours for exploration. Do not run work only to consume time, but do not treat a submit-ready candidate as completion. Preserve each trustworthy candidate as a fallback and continue scientifically meaningful exploration while the remaining budget can support experiments whose results can be completed and interpreted.

Before submitting, check the remaining budget and the plausible directions that have not yet been tested. A candidate being better than the current reference, loadable, reproducible, or artifact-valid establishes that it is a fallback; none of those facts alone establishes that exploration is complete. The default action when substantial usable budget remains is to continue exploring, analyzing, or validating.

Early submission is appropriate only when no further meaningful experiment can be completed and interpreted within the remaining budget. Do not submit merely because the current candidate is good enough or has passed its validation checks.

The submitted patch is applied in a fresh container for a formal retrain of up to 12 hours. Formal retraining starts from the fixed policy, regenerates boards, and does not reuse exploration rollouts or checkpoints.

Your submission must encode a long-running recipe designed to make meaningful use of the formal training budget. It must not normally terminate early only because of a short fixed step or epoch limit.

Your formal recipe may decide when and how often to save complete and loadable checkpoints. Save each checkpoint under `/out/checkpoints/checkpoint-<progress>/`, where `<progress>` is numeric and increases with training or construction progress.

If more than three valid checkpoints are produced, only the three with the greatest `<progress>` values will be accepted. Every accepted checkpoint will be evaluated independently, and the run's official result is the best valid final score among them. The harness handles final artifact collection and final evaluation.

Only a merged, loadable Hugging Face model is a checkpoint; raw FSDP shards are not.

## Evaluation boundary

The exploration metric is `public_four_bank_solve_rate`; the final metric is `held_out_512_board_solve_rate`. Higher is better for both. The public banks and held-out boards use different fixed environment-seed protocols, so compare each metric only with results from the same tier.

The policy start, frozen score-time Sokoban engine, action decoding, evaluation behavior, and final seeds are fixed. Candidates may change training-board construction, curriculum, rollout collection, reward shaping, objectives, and on- or off-policy updates using only information available in the training container. Formal scoring runs outside the submitted workspace. Do not import external boards, demonstrations, trajectories, or weights, reconstruct or train on final seeds, or implement an evaluation-specific lookup.

Training is stochastic at both board and policy levels. Preserve board identities and per-board outcomes, and do not treat one training seed as a complete noise estimate. A valid artifact is a merged, loadable Hugging Face checkpoint; trainer shards alone are not a result.

## Shipped solution reference

The fixed policy and the current shipped solution have the following B300 reference results:

| Measurement | Result |
|---|---:|
| Fixed policy start, final solve rate | `60/512 = 0.117188` |
| Current shipped solution, final solve rate | `87/512 = 0.169922` |
| Difference from the fixed start | `+27/512 = +0.052734` |
| Training time | `2746.19 s` |
| Final scoring time | `339.15 s` |
| Peak GPU memory during final scoring | `247,684 MiB` |

The memory number is the final-scoring peak. Training is stochastic, so report exact solved-board counts and judge a small claimed improvement against the available uncertainty before deciding what to test next.

## Work surface

Read `/workspace/run.sh`, training-board generation, rollout or data collection, advantage and reward computation, loss reduction, optimizer, checkpoint merge, and environment integration. Everything under `/workspace` is editable, including curriculum, on- or off-policy objectives, filtering, reward shaping, batching, optimization, schedule, and merge logic. These examples are illustrative, not exhaustive; they do not restrict any other change within the fixed task boundaries.

The candidate need not preserve GRPO, on-policy sampling, or the shipped training environment behavior. Formal replay must start from the fixed policy, use no external or hidden-final data, and export a merged checkpoint scored by the frozen Sokoban evaluator. Systems gains are useful only when the resulting checkpoint is evaluated under that unchanged final protocol.

## Running experiments

Give every attempt its own output tree:

```bash
OUTPUT_DIR=/out/probe-name bash /workspace/run.sh
/opt/harness/fast_eval.sh /out/probe-name/checkpoints
/opt/harness/timer.sh
```

Preserve board-bank identities, trajectory lengths, action and reward distributions, filter statistics, update timing, throughput, peak memory, trainer state, merged-checkpoint hash, evaluator payload, and failures. Training and evaluation share the GPU lock. Stop a failed candidate on environment, merge, or load failure, non-finite loss, rollout collapse, action collapse, or repeated solve-rate regression. Stopping one candidate does not by itself end exploration.

## Formal replay

Formal replay applies `candidate.patch` to a fresh `/workspace`, regenerates boards, and invokes exactly:

```bash
bash /workspace/run.sh
```

It does not inherit exploration rollouts, checkpoints, Ray state, caches, output directories, or shell exports.

## Submission

A smoke or startup check proves only that the code can begin; it is not performance evidence.

Before ending exploration, wait for every training, evaluation, and background command and read its result, or stop it explicitly and record why. Preserve the best trustworthy candidate as a fallback while exploring other directions.

Before the final action, verify that the final source starts cleanly and can merge its checkpoint.

Before submitting, verify that the patch encodes the long formal recipe and checkpoint-saving policy described above.

When no further meaningful experiment can be completed and interpreted within the remaining budget, verify the final source and artifacts, then run `/opt/harness/submit.sh` as the final action. If no candidate is trustworthy, use `/opt/harness/no_candidate.sh "reason"`. Deadline capture is recovery only and is not a normal submission path.
\end{lstlisting}

